%% ===========================================================================
%%  研究ノート・講義ノート用テンプレート（堀田グループ）
%%
%%  エンジンは LuaLaTeX です。このディレクトリで
%%      latexmk        PDF を作る
%%      latexmk -pvc   保存するたびに作り直す（書きながら使うならこれ）
%%
%%  ノートは「あとで自分が読み返して分かること」が目的なので、
%%  日付と、どこまで確かめたのかを残しておくと後で助かります。
%% ===========================================================================

\documentclass[a4paper,11pt]{ltjsarticle}

%% 共通プリアンブル（フォント・数式・図・単位・マクロ）。
%% 中身は achilles.sty を見てください。全テンプレートで共通です。
\usepackage[nogeometry]{achilles}
\usepackage[margin=22truemm]{geometry}

%% newtxmath が \openbox を定義済みなので、そのままだと amsthm と衝突する。
%% 証明終わりの記号は amsthm 側の定義を使うので、先に消しておく。
\let\openbox\relax
\usepackage{amsthm}
\usepackage{tcolorbox}
\tcbuselibrary{skins,breakable}

%% 式番号は節ごとに 1.1, 1.2 と振る
\numberwithin{equation}{section}

%% --- 定理・定義などの環境 -------------------------------------------------
\theoremstyle{definition}
\newtheorem{teigi}{定義}[section]
\newtheorem{teiri}[teigi]{定理}
\newtheorem{rei}[teigi]{例}
\theoremstyle{remark}
\newtheorem*{chui}{注意}

%% --- 覚え書き用の箱 -------------------------------------------------------
%% \begin{memo}{見出し} ... \end{memo}
%% 「なぜこうしたのか」「ここで詰まった」を書き残すのに使う。
\newtcolorbox{memo}[1][]{breakable,colback=yellow!6,colframe=yellow!50!black,
  fonttitle=\bfseries\small,title={#1},left=5pt,right=5pt}

%% --- 未確認の事項 ---------------------------------------------------------
%% achilles.sty の \todo{} と同じく、消し忘れを防ぐため目立たせている。
%% あとで grep して回収する。
\newcommand{\yokakunin}[1]{\textcolor{orange}{\textbf{[要確認: #1]}}}

\title{ノートの表題}
\author{堀田 英之}
\date{\today}

\begin{document}
\maketitle

%% 節が増えてきたら目次を出す。要らなければこの行を消す。
\tableofcontents

\section{はじめに}

このノートで何を確かめたいのかを最初に一行で書いておく。
半年後の自分は、何を考えていたか覚えていない。

\begin{memo}[この日の目的]
  対流層の熱輸送で、混合距離理論の予言と数値計算がどこでずれるのかを整理する。
\end{memo}

\section{式の展開}

\begin{teigi}[エントロピー勾配]
  超断熱度を
  \begin{align}
    \delta \equiv \nabla - \nabla_\mathrm{ad}
  \end{align}
  で定義する。
\end{teigi}

対流によるエンタルピー流束は
\begin{align}
  F_\mathrm{conv} = \rho c_p \overline{u_r T'}
  \label{eq:fconv}
\end{align}
と書ける。式\eqref{eq:fconv}の右辺を混合距離理論で評価すると
\begin{align}
  F_\mathrm{conv} \sim \rho u_r^3
\end{align}
となる。

\begin{chui}
  ここでは密度成層を無視した。実際の太陽では
  $\rho$ が対流層内で 6 桁変わるので、この近似は表面近くで破綻する。
\end{chui}

\yokakunin{係数の $1/2$ がどこから来るか、Kippenhahn の教科書で確認する}

\section{数値の確認}

\begin{rei}[太陽の対流速度]
  $F_\odot = \qty{6.3e7}{\watt\per\meter\squared}$、
  $\rho = \qty{0.2}{\kg\per\meter\cubed}$ とすると
  $u_r \sim (F_\odot/\rho)^{1/3} \simeq \qty{700}{\meter\per\second}$ となり、
  観測される粒状斑の速度と同程度である。
\end{rei}

\section{図を貼る}

%% 画像はこのディレクトリか fig/ に置く。
%% LuaLaTeX が読めるのは PDF・PNG・JPEG。EPS は先に PDF へ変換しておく。
%\begin{figure}[tb]
%  \centering
%  \includegraphics[width=0.7\linewidth]{foo.pdf}
%  \caption{図の説明。}
%  \label{fig:foo}
%\end{figure}

\end{document}
